\section{Patching the selected local lattices into one global lattice}
\label{sec:lattice-patching}

\status{LP1--LP5 prove the finite-place patching and its exact application
to the selected source lattice. The global geometric datum, the local
comparison and its commensurability bounds are supplied. This does not
construct the geometric point or the stronger numerical record
\texttt{NumAt}.}

Write $R=\mathcal O_K$ for the integer ring of a number field, and let
$V$ be a $K$-vector space. For an $R$-submodule $M\subset V$ and a
nonzero prime ideal $\mathfrak p$, identify its localization with
\[
 M_{\mathfrak p}=\{x\in V:\text{some }s\in R\setminus\mathfrak p
                                      \text{ satisfies }sx\in M\}.
\]
All intersections below are inside $V$. A module $A_{\mathfrak p}$ is
called local at $\mathfrak p$ when $sx\in A_{\mathfrak p}$ and
$s\notin\mathfrak p$ imply $x\in A_{\mathfrak p}$. This is closure
under the units inverted in $R_{\mathfrak p}$; it does not mean
saturation with respect to every nonzero element of $R$.

\paragraph{LP1: clear every denominator except at one chosen prime.}
For $0\ne c\in R$ and a fixed $\mathfrak p$, there is
$t\notin\mathfrak p$ such that, for every $\mathfrak q\ne\mathfrak p$,
some $s\notin\mathfrak q$ satisfies $c\mid st$.
Indeed factor
\[
                 (c)=\mathfrak p^a\mathfrak a,
                 \qquad\mathfrak a+\mathfrak p=R.
\]
Choose $t\in\mathfrak a\setminus\mathfrak p$. If
$\mathfrak q\ne\mathfrak p$, then
$\mathfrak p^a\not\subset\mathfrak q$: for $a>0$ the contrary
would imply $\mathfrak p\subset\mathfrak q$ and hence equality;
for $a=0$ it would say $R\subset\mathfrak q$.
Choose $s\in\mathfrak p^a\setminus\mathfrak q$. Then
$st\in(c)$, as required. No principal generator for
$\mathfrak p$ is chosen.

\paragraph{LP2: localization commutes with this particular intersection.}
Suppose $B_{\mathfrak q}\subset V$ is local at $\mathfrak q$ for
every nonzero prime, and that every $x\in V$ has a common nonzero
denominator:
\[
 \forall x\in V\quad\exists,0\ne c\in R\quad
                   cx\in B_{\mathfrak q}\ \text{for every }\mathfrak q.
                                                               \tag{LP.1}
\]
Then
\[
 N:=\bigcap_{\mathfrak q}B_{\mathfrak q}
       \quad\Longrightarrow\quad N_{\mathfrak p}=B_{\mathfrak p}.
                                                               \tag{LP.2}
\]
One inclusion follows from the local property. For the other, take
$x\in B_{\mathfrak p}$ and choose $c$ in (LP.1). LP1 gives
$t\notin\mathfrak p$. At $\mathfrak p$ we have
$tx\in B_{\mathfrak p}$ directly. At every other
$\mathfrak q$, write $st=cd$ with $s\notin\mathfrak q$.
Then $s(tx)=d(cx)\in B_{\mathfrak q}$, and the local property
gives $tx\in B_{\mathfrak q}$. Thus $tx\in N$, proving (LP.2).
The denominator hypothesis is the reason this infinite intersection
may be localized; no unrestricted commutation rule is being used.

\paragraph{LP3: finite modification of a full reference module.}
Let $S$ be a finite set of primes and suppose that $M$ is full in the
explicit sense that every $x\in V$ has a nonzero integer multiple in
$M$. For $\mathfrak p\in S$, prescribe a local module
$A_{\mathfrak p}$ and a positive integer $n_{\mathfrak p}$ with
\[
 n_{\mathfrak p}M\subset A_{\mathfrak p},\qquad
 n_{\mathfrak p}A_{\mathfrak p}\subset M_{\mathfrak p}.
                                                               \tag{LP.3}
\]
Set $B_{\mathfrak p}=A_{\mathfrak p}$ on $S$ and
$B_{\mathfrak p}=M_{\mathfrak p}$ off $S$. If $mx\in M$ and
$d=\prod_{\mathfrak p\in S}n_{\mathfrak p}$, then $dmx$ belongs
to every $B_{\mathfrak p}$. Hence LP2 constructs the exact patch
\[
 \boxed{\quad
 N_{\mathfrak p}=A_{\mathfrak p}\ (\mathfrak p\in S),\qquad
 N_{\mathfrak p}=M_{\mathfrak p}\ (\mathfrak p\notin S).
 \quad}                                                       \tag{LP.4}
\]
There is also the useful global bound
\[
                        dM\subset N\subset d^{-1}M.            \tag{LP.5}
\]
For the right inclusion, $dx\in M_{\mathfrak p}$ at every prime
when $x\in N$. To justify the resulting global membership, observe
that $M=\bigcap_{\mathfrak p}M_{\mathfrak p}$ for any submodule:
if $x\notin M$, the proper ideal $\{r:rx\in M\}$ is contained
in a maximal ideal, which contradicts $x$ belonging to that
localization. In the present full situation this ideal contains a
nonzero integer, so the maximal ideal is nonzero. The same argument
proves uniqueness of $N$ from (LP.4).
If $M$ is a finite full lattice, (LP.5) and Noetherianity make $N$
a finite full lattice too. The intersection proof of existence
itself does not require finite generation.

\paragraph{LP4: the literal served-prime and type-prime patch.}
Fix a supplied word $u$ and put
$M=M_{\rm ram}(u)\subset V_u$, the retained reference lattice of
PL4. This is the lattice called \texttt{DV.MramL}; it is not the
whole unretained source lattice. Let $T_u$ be its own type prime.
For a good place $w$ at a served rational prime, prescribe
\[
 A_w=\{x\in V_u:\operatorname{cmp}_{u,w}(\nu_u x)
                                    \in\mathcal N_{u,w}\},       \tag{LP.6}
\]
where $\mathcal N_{u,w}$ is the actual selected local module
\texttt{NP}, and served means the literal set \texttt{VV} is
nonempty. At a place above $T_u$, prescribe instead
\[
                  A_w=\{x:\nu_u x\in L^\lambda_{u,w}\}.         \tag{LP.7}
\]
The valuation compatibility of the comparison identifies every
scalar integral at $w$ with a scalar of the coefficient DVR.
The modules in (LP.6)--(LP.7) are therefore closed under these
scalars, hence local at $w$. This proves the local saturation
condition used in LP2 directly, including for denominators outside
$w$.

The supplied served-place conditions say
$A_w\subset M_w$ and $n_wM\subset A_w$ for some $n_w>0$.
They imply both clauses of (LP.3). The supplied type-place conditions
say exactly $n_wM\subset A_w$ and $n_wA_w\subset M_w$.
Let $S$ consist of these served places and the places above $T_u$.
It is finite: a served prime satisfies $p<X'$, and a number field
has finitely many places above each rational prime. Moreover the
good-place condition excludes $p\mid T_u$, so (LP.6) and (LP.7)
never compete. LP3 now gives one $N_u$ satisfying both prescriptions,
with
\[
 \begin{array}{ll}
 (N_u)_w\subset M_w,&p_w\ne T_u,\\
 (N_u)_w=M_w,&p_w\ne T_u\text{ and }\texttt{VV}(p_w)=\varnothing.
 \end{array}                                                   \tag{LP.8}
\]
No containment in $M_w$ is imposed at the type prime. In particular
(LP.8) cannot be strengthened to $N_u\subset M$ without an
additional type-place assumption. Nor is $N_u$ generally saturated:
already the prescription $A_{(2)}=2\mathbb Z_{(2)}$, with reference
$\mathbb Z$ and no other modifications, gives $N=2\mathbb Z$.

\paragraph{LP5: install the patch without losing the valuation scale.}
Once the canonical refilled size/comparison datum is supplied, replace
its \texttt{NuLat} field by the modules $N_u$ above. The source
\texttt{setLat} sets
\[
 \operatorname{lenMN}(u,w)
  =\frac{\operatorname{toNat}\bigl(\operatorname{lenLoc}_w
                                      (N_u,M_{\rm ram}(u))\bigr)}{e_w},
              \qquad e_w>0.                                  \tag{LP.9}
\]
Here \texttt{lenLoc} is the source's extended-natural local length
and \texttt{toNat} its conversion; (LP.9) does not assert finiteness
at an arbitrary place. Whenever this local length is a finite integer
$\ell$, substitution gives the exact required identity
\[
                   \ell=e_w\operatorname{lenMN}(u,w).           \tag{LP.10}
\]
At a non-type unserved place the two localized modules agree by
(LP.8), so the length is zero. At served places the actual integral
comparison is exactly (LP.6), not merely an abstract isomorphism.
All other fields of the size datum are retained. Consequently its
scale, selected degrees, rank and host bounds, and canonical orbit
selector are unchanged; the instance is transported by keeping its
same field values. These are the lattice and length portions of
\texttt{exists\_ledPart}.

\paragraph{Exact boundary of the point construction.}
LP4 uses the local containment and commensurability assertions in
\texttt{RecOK.atT} and \texttt{LedOK}, and LP5 uses a supplied
canonical refilled datum. The source derives the latter through
\texttt{exists\_refill}, which installs finite support and the gain,
error and deficit fields. Neither the construction of the underlying
\texttt{Lay}/\texttt{RecOK}/\texttt{LedOK} data nor that complete
refill theorem is claimed here. The later passage
\texttt{RecExP}\,$\Rightarrow$\,\texttt{ObjCP} also needs the level
prime and size constructors; \texttt{NumAt} additionally needs its
explicit parameter and block-length inequalities. These remain
separate targets. What has been removed is the need to assume that
the already specified local lattices can be simultaneously patched
into a global one.
